\subsection{Element Model}\label{sec:genwarp}

The governing equations for the standard Type II family of beam theories is briefly recalled. The objective is to obtain a model on a curve \(\mathcal{C}\) that represents finite deformations of a 3D body \(\Body\).
In \cref{sec:std-continuum}, the body is introduced in the continuum setting. \Cref{sec:std-kinematics} introduces the kinematic representation of this body in terms of fields on the curve \(\mathcal{C}\), from which the beam deformations \(\bm{e}\) naturally precipitate and conjugate generalized stress \(\bm{s}\) are implied. \Cref{sec:std-equilibrium} derives equilibrium equations for these resultants \(\bm{s}\). \Cref{sec:std-section} derives the section constitutive relation \(\bm{e}\rightarrow\bm{s}\) by completing the virtual work identity.

\subsubsection{Continuum}\label{sec:std-continuum}
In the continuum setting the model is a body
\(
\Body=(0,L)\times\Section
\),
where \(\Section\subset\mathbb{R}^3\) is a bounded cross-section
spanned by basis vectors \(\{\FrameBasis_y,\FrameBasis_z\}\).% in \(\mathbb{R}^3\). 
Together with the axial unit vector \(\FrameBasis\), the vectors \(\{\FrameBasis,\,\FrameBasis_y,\,\FrameBasis_z\}\) form an orthonormal basis and the 3D identity is:
\[
\mathbf{1}=\FrameBasis\otimes\FrameBasis+\oneS
\qquad\text{with}\qquad
\oneS=\FrameBasis_{\beta}\otimes\FrameBasis_{\beta}
\]
A point in \(\Body\) has coordinates \(\reflenx\,\FrameBasis + \bm{r}\) in the reference configuration, where \(\bm{r} \in\Section_x\) is the coordinate of a point in the section at \(\reflenx\) and \(\bm{r}\cdot\FrameBasis = 0\).
The boundary of \(\Body\) consists of the lateral boundary \(\BodyLateralSurface=\SectionBoundary\times(0,L)\) and the surfaces \(\Section_0,\,\Section_L\) of the sections at the ends \(\reflenx \in \{0,L\}\).  

\begin{remark}
The following developments are made without appealing to a coordinate system (we only assume right-handedness), and thus no commitments are made to a particular coordinate convention. However, such conventions can aid the interpretation of equations and are required for implementation. 
Two common conventions for assigning coordinates \(x,y,z\) are:
\begin{itemize}
    \item Take the first (\(x\)) coordinate to coincide with the reference longitudinal axis \(\FrameBasis\). This furnishes components on the basis \(\{\bm{\partial}_x,\bm{\partial}_y,\bm{\partial}_z\}\)
    \[
    \FrameBasis \equiv \bm{\partial}_x = \begin{pmatrix}
        1 \\ 0 \\ 0
    \end{pmatrix},
    \quad
    \PlanePosition = \begin{pmatrix}
    0 \\ y \\ z
    \end{pmatrix}
    \qquad\text{and}\qquad 
    \FrameBasis^{\times} = \begin{bmatrix}
        0 & 0 & 0 \\
        0 & 0 & -1 \\
        0 & 1 & 0
    \end{bmatrix}
    \]
    \item Take the last (\(z\)) coordinate to coincide with \(\FrameBasis\). This furnishes components on the basis \(\{\bm{\partial}_x,\bm{\partial}_y,\bm{\partial}_z\}\)
    
    \[
    \FrameBasis \equiv \bm{\partial}_z = \begin{pmatrix}
        0 \\ 0 \\ 1
    \end{pmatrix},
    \quad
    \PlanePosition = \begin{pmatrix}
    x \\ y \\ 0
    \end{pmatrix}
    \qquad\text{and}\qquad 
    \FrameBasis^{\times} = \begin{bmatrix}
        0 & -1 & 0 \\
        1 &  0 & 0 \\
        0 & 0 & 0
    \end{bmatrix}
    \]
\end{itemize}
\end{remark}

\begin{remark}
The 3D out-of-plane reference director \(\FrameBasis\) is used to form cross products with vectors in the plane \(\Section\). 
For example, given a function \(u: S \rightarrow \mathbb{R}\) with \(u' = 0\), the notation \(\FrameBasis \times \nabla u\) may be interpreted in coordinates \(\bm{\partial}_x \coloneq \FrameBasis\) either as 
\[
\begin{bmatrix}
0 & 0 & 0 \\
0 & 0 & -1 \\
0 & 1 & 0
\end{bmatrix}
\begin{pmatrix}
0 \\ u_{,y} \\ u_{,z}
\end{pmatrix}
\qquad\text{or}\qquad 
\begin{bmatrix}
0 & -1 \\
1 & 0
\end{bmatrix}
\begin{pmatrix}
u_{,y} \\ u_{,z}
\end{pmatrix}
\]
\end{remark}



\subsubsection{Kinematics}\label{sec:std-kinematics}

The deformation is represented by the deflection \(\bm{x}(\reflenx)\in\mathbb{R}^3\), rotation \(\FrameRotation(\reflenx)\in\mathrm{SO}(3)\), and warping fields \(\bm{\alpha}(\reflenx)\in\mathbb{R}^{n_w}\):
\begin{SaveEquation}{eq:finite-alignment-field}
\hat{\bm{x}}(\reflenx, \PlanePosition) 
= \bm{x}(\reflenx) + \FrameRotation \bigl(\PlanePosition 
+\FrameBasis\otimes\WarpAxialShapes(\bm{r})\bm{\alpha}(\reflenx)
% +\WarpTotalShapes(\bm{r})\bm{\alpha}(\reflenx)
\bigr)
\end{SaveEquation}
where \(\WarpAxialShapes:\Section\rightarrow\mathbb{R}^{n_w}\) collects \(n_w\) cross-sectional warping shapes. 
% The representation \eqref{eq:finite-alignment-field} is exact in the limit \(n_w\rightarrow\infty\). 
%
\iffalse
The warping shapes \(\bm{\Phi}\) may be decomposed into axial and in-plane parts:
\begin{equation}
\WarpTotalShapes = \FrameBasis\otimes\WarpAxialShapes + \WarpPlaneShapes,
\qquad
\WarpAxialShapes\coloneq \FrameBasis\cdot\bm{\Phi}\in\mathbb{R}^{1\times n_w},
\label{eq:decom-warp-shapes}
\end{equation}
where \(\WarpAxialShapes\) collects the axial components and \(\WarpPlaneShapes\) the in-plane components of the warping shapes. 
Correspondingly, the warping contribution to the deformation \eqref{eq:finite-alignment-field} decomposes into a scalar axial warping field and a vector in-plane warping field:
\begin{equation*}
\AxialWarp\coloneq \WarpAxialShapes\cdot\bm{\alpha} = \FrameBasis\cdot\bm{\Phi}\bm{\alpha}\in\mathbb{R}, 
\qquad\text{and}\qquad
\PlaneWarp\coloneq \mathbf{1}_{\Section}\,\bm{\Phi}\bm{\alpha} = \WarpPlaneShapes\,\bm{\alpha}.
\end{equation*}

The axial warping shapes \(\WarpAxialShapes\) enrich the axial and shear strain components---those coupling to \(\FrameBasis\)---and thus directly participate in the beam's virtual work. 
The in-plane warping \(\PlaneWarp\), by contrast, contributes only membrane-type (\(\alpha\beta\)) components to the material strain tensor. 
As will be shown in \cref{sec:sv-stress}, these components are energetically orthogonal to the admissible stress field under the Saint-Venant stress assumption, and are therefore invisible to the beam-level virtual work principle. 
This observation motivates restricting the \emph{standard} warping model to axial shapes \(\WarpAxialShapes\), while the effect of in-plane warping can be recovered through the enhanced strain framework of \cref{sec:enhan-section}.
\fi
%
% \subsubsection{Deformation Gradient}\label{sec:def-gradient}
The tangent map \(\bm{F}\coloneq T\,\hat{\chi}\) of the finite deformation \eqref{eq:finite-alignment-field} is:
\begin{equation}
\bm{F}
=\FrameRotation \left(\mathbf{1}+\bm{\mu}\otimes\FrameBasis
~+~\FrameBasis\otimes\nabla\AxialWarp
% ~+~\nabla\PlaneWarp
\right),
\qquad\text{with}\qquad
\bm{\mu}(\bm{r})
=\ShearStrain+\CurveStrain\times\bigl(\bm{r}+\FrameBasis\otimes\WarpAxialShapes\,\bm{\alpha}\bigr)
% +\FrameBasis\otimes\WarpAxialShapes\,\bm{\alpha}'
\label{eq:F-Phi}
\end{equation}
where \(\AxialWarp := \WarpAxialShapes\cdot\bm{\alpha}\) is the magnitude of axial warping and 
\(\ShearStrain\) and \(\CurveStrain\) are the beam strain measures:
\begin{SaveEquation}{eq:def-gam-kap}
  \begin{aligned}
    \bm{\gamma} &= \FrameRotation^{\top} \bm{x}^\prime - \FrameBasis  \\
    &= \epsilon \FrameBasis + \gamma_\beta \FrameBasis_\beta 
  \end{aligned}
  \qquad\text{ and }\qquad
  \begin{aligned}
    \bm{\kappa} &= \left(\FrameRotation^{\top} \FrameRotation^\prime \right)^{\vee} \\
    &= \Twist \FrameBasis + \kappa_\beta  \FrameBasis_{\beta }
  \end{aligned}
\end{SaveEquation}
The vector \(\ShearStrain\) comprises the axial strain \(\epsilon\) and the transverse shears \(\gamma_\beta\), and \(\CurveStrain\) comprises the twist \(\Twist\) and curvatures \(\kappa_\beta\).
% \eqref{eq:def-gam-kap} and \(\nabla(\cdot)\) acts on section coordinates \(\bm{r}\) only.
%
Note that for $\bm{\AxialWarp} \coloneq (\bm{\alpha}(\reflenx)\cdot\WarpAxialShapes(\bm{r}))\FrameBasis$
$$
\nabla \bm{\AxialWarp} = \nabla (\bm{\alpha}\cdot\WarpAxialShapes\FrameBasis) = \FrameBasis\otimes\nabla(\bm{\alpha}\cdot\WarpAxialShapes) = \FrameBasis\otimes\nabla\AxialWarp
$$
\begin{Details}
$$
\nabla \AxialWarp = (\nabla\bm{\alpha})^{\top}\WarpAxialShapes  +(\nabla\WarpAxialShapes)^{\top}\bm{\alpha} 
$$
Using $\nabla \bm{\alpha} = \bm{\alpha}'\otimes\FrameBasis$
$$
\nabla \AxialWarp = (\bm{\alpha}^{\prime}\cdot\WarpAxialShapes) \FrameBasis +(\nabla\WarpAxialShapes)^{\top}\bm{\alpha}
$$
and rearranging
$$
\nabla \AxialWarp = (\FrameBasis\otimes\WarpAxialShapes) \bm{\alpha}' + (\nabla\WarpAxialShapes)^{\top}\bm{\alpha}
$$

$$
(\mathbf{a}\stackrel{\mathrm{s}}{\otimes}\FrameBasis)\FrameBasis
=\frac{1}{2} \left(\mathbf{a} + (\mathbf{a}\cdot\FrameBasis)\FrameBasis\right)
$$

$$
\delta \left((\mathbf{a}\stackrel{\mathrm{s}}{\otimes}\FrameBasis)\FrameBasis\right) = \frac{1}{2}\left(\delta \mathbf{a} + \delta a_1 \FrameBasis\right)
$$
\end{Details}
%
% \subsubsection{Material Strain Tensor}\label{sec:material-strain}

The Green-Lagrange strain tensor \(\hat{\StrainTensor}\coloneq\tfrac12(\bm{F}^{\top}\bm{F}-\mathbf{1})\) with \eqref{eq:F-Phi} is:
\begin{equation}
\begin{aligned}
\hat{\StrainTensor} = \bm{\mu}
\stackrel{\mathrm{s}}{\otimes}
\FrameBasis 
+\nabla \AxialWarp  \stackrel{\mathrm{s}}{\otimes} \FrameBasis
+({\bm{\mu}}\cdot\FrameBasis)\,\nabla \AxialWarp  \stackrel{\mathrm{s}}{\otimes} \FrameBasis
+ \frac{1}{2}\left(
  \bm{\mu}\cdot\bm{\mu}\,\FrameBasis\otimes\FrameBasis
+ \nabla \AxialWarp \otimes\nabla \AxialWarp
\right)
% \hat{\StrainTensor}
% &=~\bm{\mu}\stackrel{\mathrm{s}}{\otimes}\FrameBasis
% ~+~\nabla\AxialWarp\stackrel{\mathrm{s}}{\otimes}\FrameBasis
% ~+\stackrel{\rm s}{\nabla}\PlaneWarp
% \\
% &\quad
% +~(\bm{\mu}\cdot\FrameBasis)\,\nabla\AxialWarp\stackrel{\mathrm{s}}{\otimes}\FrameBasis
% ~+~\tfrac12\Big(
% \bm{\mu}\cdot\bm{\mu}\,\FrameBasis\otimes\FrameBasis
% + \nabla\AxialWarp\otimes\nabla\AxialWarp
% + (\nabla\PlaneWarp)^{\top}\nabla\PlaneWarp
% \Big),
\end{aligned}
\label{eq:exact-gl-strain-Phi}
\end{equation}
where the symmetric bun product \(\bm{a}\stackrel{\mathrm{s}}{\otimes}\bm{b}\) is defined for vectors \(\bm{a}\), \(\bm{b}\), and \(\bm{c}\) by
\begin{equation*}
    (\bm{a}\stackrel{\mathrm{s}}{\otimes}\bm{b} )\bm{c} \coloneq \tfrac{1}{2} \bigl(
      \bm{a}\,(\bm{b}\cdot\bm{c}) + \bm{b}\,(\bm{a}\cdot\bm{c})
    \bigr).
\end{equation*}

% The exact strain tensor \eqref{eq:exact-gl-strain-Phi} admits a natural decomposition:
% \begin{equation}
% \hat{\StrainTensor}
% = \hat{\StrainVector}\stackrel{\mathrm{s}}{\otimes}\FrameBasis
% ~+~\stackrel{\rm s}{\nabla}\PlaneWarp
% ~+~\text{h.o.t.}
% \label{eq:strain-decomposition}
% \end{equation}
% where the first term contains the components coupling to \(\FrameBasis\), and the second term \(\stackrel{\rm s}{\nabla}\PlaneWarp\) has only in-plane (\(\alpha\beta\)) components. 
% \begin{remark}
% \label{rem:empty-strains}
% From \eqref{eq:exact-gl-strain-Phi} it is evident that a displacement field with the form \eqref{eq:finite-alignment-field} is incapable of producing strains orthogonal to the axis of the rod \(\FrameBasis\). 
% That is, \(\FrameBasis_{\alpha} \cdot \hat{\StrainTensor} \FrameBasis_{\beta} = 0\) and the  strain \cref{eq:exact-gl-strain-Phi} is represented by a vector \(\hat{\StrainVector} \in \mathbb{R}^3\) with the form:
% \begin{equation*}
%     \hat{\StrainTensor} = \hat{\StrainVector}\symbun\FrameBasis
% \end{equation*}
% This is not generally representative of slender bodies and may be rectified by introducing an ``enhanced'' (incompatible) field \(\PlaneStrainTensor\) that is obtained by enforcing the stress constraints \(\FrameBasis_{\alpha} \cdot \hat{\StrainTensor} \FrameBasis_{\beta} = 0\). 
% % Although this is not generally discussed in the enhanced/assumed strain context, the procedure is common practice and will not be extensively developed here. See, eg, \cite{klinkel2002using}.
% \end{remark}

\begin{remark}
The linearization of \eqref{eq:finite-alignment-field} is:
\begin{equation}
\hat{\bm x}(\reflenx, \bm{r})=(\reflenx\FrameBasis+\bm{r})
+\bm{u} + \bm\theta(\reflenx)\times\bm{r}
% +\WarpTotalShapes(\bm{r})\,\bm\alpha(\reflenx)
+\FrameBasis\otimes\WarpAxialShapes(\bm{r})\,\bm\alpha(\reflenx)
\label{eq:linear-alignment-field}
\end{equation}
and the linearized displacement \(\hat{\bm{u}} = \hat{\bm{x}} - (\reflenx\FrameBasis + \bm{r})\) is
\begin{SaveEquation}{eq:trace-lift-linear}
\hat{\bm u}(\reflenx, \bm{r})=\bm{u}(\reflenx)
+\bm\theta(\reflenx)\times\bm{r}
% +\WarpTotalShapes(\bm{r})\,\bm{\alpha}(\reflenx)
+\FrameBasis\otimes\WarpAxialShapes(\bm{r})\,\bm\alpha(\reflenx)
\end{SaveEquation}
In terms of \(\bm{u}\triangleq\bm{x}- \reflenx\,\FrameBasis\) the linearized beam strains are:
\begin{SaveEquation}{eq:gam-kap-linear}
\bm{\gamma}=\bm{u}'+\FrameBasis\times\bm\theta,\qquad
\bm{\kappa}=\bm\theta'.
\end{SaveEquation}
\end{remark}

\begin{remark}
The nonlinearities in the strain tensor \eqref{eq:exact-gl-strain-Phi} are of two distinct origins:
\begin{itemize}
    \item \emph{Element-level nonlinearity}: encoded in the beam strains \(\ShearStrain\), \(\CurveStrain\) and resolved by the finite element discretization of the beam axis.
    This includes the finite rotation \(\FrameRotation\) and the strain measures \eqref{eq:def-gam-kap}.
    \item \emph{Section-level nonlinearity}: arising from products of section-coordinate functions (\(\bm{r}\), \(\nabla\AxialWarp\), etc.) with beam strain components, and resolved \emph{locally} at each cross section. The Wagner term \(\tfrac{1}{2}\Twist^2\,\bm{r}\cdot\bm{r}\,\FrameBasis\) is the prototypical example: it is quadratic in the twist \(\Twist\) and quadratic in the section coordinate \(\bm{r}\), but does not introduce coupling between different points along the beam axis.
\end{itemize}
This separation allows the section constitutive model to be developed independently of the element formulation. 
All strain representations derived below may be employed in connection with either the finite deformation beam model \eqref{eq:finite-alignment-field} or its linearization \eqref{eq:trace-lift-linear}.
\end{remark}



\subsubsection{Equilibrium Equations}\label{sec:std-equilibrium}
% Introduce e and s
In \cref{sec:std-kinematics} the measures \(\ShearStrain\), \(\CurveStrain\), \(\bm{\alpha}\) and \(\bm{\alpha}'\) arise as natural strain measures on the curve \(\mathcal{C}\). 
Thus the state of stress and strain along \(\reflenx\) is modeled by:
\begin{SaveEquation}{eq:cosserat-strain} 
  \bm{s} \coloneq \begin{pmatrix}
  \bm{n} \\ 
  \bm{m} \\
  \bm{w} \\
  \bm{v}
  \end{pmatrix}
  \quad\text{ and }\quad 
  \bm{e} \coloneq \begin{pmatrix}
  \bm{\gamma}\\ 
  \bm{\kappa} \\
  \bm{\alpha}' \\
  \bm{\alpha}
  \end{pmatrix}
  % \label{eq:cosserat-strain}
\end{SaveEquation}
where \(\bm{n}\),\(\bm{m}\), ...

The governing equations of the beam theory are obtained by requiring that the virtual work identity
\begin{equation}
\left<\TraceTotalEnergyStress,\,\delta\FrameStrain\right>_{\FrameBody}=\left<\StressTensor ,\,\delta\StrainTensor\right>_{\hat{\FrameBody}}
\label{eq:seSE}
\end{equation}
hold for all admissible variations \(\delta\bm{e}\). 

The identity \eqref{eq:seSE} is enforced in two steps: integration over the cross section yields the stress resultants \eqref{eq:stress-resultants}, and subsequent integration along the beam axis yields the equilibrium equations \eqref{eq:finite-equilibrium}.


% \subsubsection{Resultants}\label{sec:resultants}

% This furnishes the equalities
% \begin{equation*}
% \left<\TraceTotalEnergyStress,\,\delta\FrameStrain\right>
% \equiv \left<\StressVector,\, \delta \hat{\StrainVector}\right>
% \equiv \left<\StressTensor ,\,\delta\StrainTensor\right>.
% \end{equation*}


% \subsubsection{Equilibrium}\label{sec:equilibrium}

Imposing the remainder of \eqref{eq:seSE} by integration along \(\reflenx\) and integrating by parts furnishes the complete system of equilibrium equations:
\begin{equation} 
\left\{\begin{array}{r}
(\FrameRotation \bm{n})^{\prime}+{\bm{f}}_n=\mathbf{0} \\
(\FrameRotation \bm{m})^{\prime}+\bm{x}^{\prime} \times (\FrameRotation \bm{n}) +{\bm{f}}_m=\mathbf{0} \\
\bm{w}^{\prime}-\bm{v} + \bm{f}_{\rm w} = \mathbf{0}
\end{array}\right. 
\label{eq:finite-equilibrium}
\end{equation}
where \(\bm{f}_n\) and \(\bm{f}_m\) are given body force and moment loadings, respectively. 


\subsubsection{Material Constitution}\label{sec:std-material}

The general state of strain has the decomposition:
\begin{SaveEquation}{eq:strain-tensor-decomp}
\hat{\StrainTensor} = \FrameBasis\stackrel{s}{\otimes}(\PointAxialStrain\,\FrameBasis 
+ \PointShearStrain) + \PlaneStrainTensor_{\alpha\beta}\FrameBasis_{\alpha}\otimes\FrameBasis_{\beta} 
\end{SaveEquation}
with \(\PointShearStrain\) a vector with zero projection on \(\FrameBasis\). 
The compatible strain \eqref{eq:exact-gl-strain-Phi} contributes to both parts: 
the terms involving \(\stackrel{\mathrm{s}}{\otimes}\FrameBasis\) furnish the 
axial and shear components, while the quadratic term 
\(\tfrac{1}{2}\nabla\AxialWarp\otimes\nabla\AxialWarp\) produces in-plane 
(\(\alpha\beta\)) components.
In addition, the physical strain field includes further in-plane contributions 
(e.g., lateral contraction) that are not captured by the compatible field 
\eqref{eq:finite-alignment-field}.

It is standard to restrict the material response to stress tensors with the Saint-Venant form:
\begin{SaveEquation}{eq:saint-venant-stress}
\StressTensor = \AxialStress\, \FrameBasis \otimes \FrameBasis+\ShearStress \otimes \FrameBasis+\FrameBasis \otimes \ShearStress
\qquad\text{with}\qquad
\begin{aligned}
\AxialStress    &\coloneq \FrameBasis\cdot \StressTensor\, \FrameBasis \\
\ShearStress &\coloneq \FrameBasis_{\beta}\otimes \FrameBasis_{\beta}\, \StressTensor\,\FrameBasis
\end{aligned}
\end{SaveEquation}
where \(\AxialStress\) is the scalar axial stress and \(\ShearStress \in \mathbb{R}^2\) the vector of shear stresses. 
Thus, the in-plane stress components \(\FrameBasis_{\alpha}\cdot\StressTensor\,\FrameBasis_{\beta}\) are assumed to vanish. 
This furnishes the stress vector \(\StressVector \in \mathbb{R}^3\):
\begin{equation}
\StressVector := \StressTensor\,\FrameBasis = \AxialStress \FrameBasis + \ShearStress.
\label{eq:def-stress-vector}
\end{equation}

Since \(\StressTensor_{\alpha\beta}=0\), neither the compatible in-plane strain 
\(\tfrac{1}{2}\nabla\AxialWarp\otimes\nabla\AxialWarp\) nor any additional 
in-plane strain \(\PlaneStrainTensor_{\alpha\beta}\) contributes to the virtual work. 
Only the strain vector \(\hat{\StrainVector}\), which carries the components 
coupling to \(\FrameBasis\), is energetically active:
\begin{equation}
\left<\StressTensor,\,\delta\hat{\bm{E}}\right>
=\left<\StressVector,\,\delta\hat{\bm{\varepsilon}}\right>.
\label{eq:sv-vw-reduction}
\end{equation}
The in-plane strains nevertheless influence the axial stress \(\AxialStress\) 
through the material law (e.g., via Poisson contraction). 
They are determined implicitly by enforcing \(\FrameBasis_{\alpha}\cdot\StressTensor\,\FrameBasis_{\beta}=0\) 
at each material point---a local condensation procedure detailed 
in \cref{sec:enhan-section} \citep{klinkel2002using}.
This identity is the basis for the section constitutive relations developed in \cref{sec:std-section}.

\iffalse
\subsubsection{Material Constitution}

The general state of strain has the decomposition:
\begin{SaveEquation}{eq:strain-tensor-decomp}
\hat{\StrainTensor} = \FrameBasis\stackrel{s}{\otimes}(\PointAxialStrain\,\FrameBasis 
+ \PointShearStrain) + \PlaneStrainTensor_{\alpha\beta}\FrameBasis_{\alpha}\otimes\FrameBasis_{\beta} 
\end{SaveEquation}
with $\PointShearStrain$ a vector with zero projection on $\FrameBasis$. 

\begin{equation}
\left<\StressTensor,\,\delta\hat{\bm{E}}\right>
=\left<\StressVector,\,\delta\hat{\bm{\varepsilon}}\right>,
\label{eq:sv-vw-reduction}
\end{equation}

% \paragraph{Stress Assumption}
The stress tensor is assumed to take the form:
\begin{SaveEquation}{eq:saint-venant-stress}
\StressTensor = \AxialStress\, \FrameBasis \otimes \FrameBasis+\ShearStress \otimes \FrameBasis+\FrameBasis \otimes \ShearStress
\qquad\text{with}\qquad
\begin{aligned}
\AxialStress    &\coloneq \FrameBasis\cdot \StressTensor\, \FrameBasis \\
\ShearStress &\coloneq \FrameBasis_{\beta}\otimes \FrameBasis_{\beta}\, \StressTensor\,\FrameBasis
\end{aligned}
\end{SaveEquation}
where \(\AxialStress\) is the scalar axial stress and \(\ShearStress \in \mathbb{R}^2\) the vector of shear stresses. 
Thus, the in-plane stress components \(\FrameBasis_{\alpha}\cdot\StressTensor\,\FrameBasis_{\beta}\) are assumed to vanish. 
This furnishes the representation for the stress vector \(\StressVector \in \mathbb{R}^3\):
\begin{equation}
\StressVector := \StressTensor\,\FrameBasis = \AxialStress \FrameBasis + \ShearStress.
\label{eq:def-stress-vector}
\end{equation}
\fi
\subsubsection{Section Constitution}\label{sec:std-section}

By the virtual work reduction \eqref{eq:sv-vw-reduction}, the section constitutive 
relation requires only the strain vector \(\hat{\StrainVector}\)---the part of 
\(\hat{\StrainTensor}\) coupling to \(\FrameBasis\). The compatible in-plane 
contributions (such as \(\tfrac{1}{2}\nabla\AxialWarp\otimes\nabla\AxialWarp\)) 
and the condensed strains \(\PlaneStrainTensor_{\alpha\beta}\) are invisible to 
the virtual work and need not appear explicitly. 
Extracting from \eqref{eq:exact-gl-strain-Phi} only the terms involving 
\(\stackrel{\mathrm{s}}{\otimes}\FrameBasis\), one obtains the strain vector 
\(\hat{\StrainVector}\) as a function of the beam strains \(\bm{e}\) and the 
section coordinates \(\bm{r}\). This expression may be truncated at various 
orders of section-level nonlinearity; two cases of interest are:

\begin{description}
\item[Moderate (Wagner) strain.]
A truncation of \eqref{eq:exact-gl-strain-Phi} which retains the dominant section-level nonlinearity is:
\begin{SaveEquation}{eq:wagner-gl-strain-pi}
%
\hat{\StrainVector}\coloneq
\bm{\gamma}
+\bm{\kappa}\times\bm{r}
+\FrameBasis\otimes\WarpAxialShapes\,\bm{\alpha}'
+\gradS (\WarpAxialShapes\cdot\bm{\alpha})
+\tfrac{1}{2}\Twist^2\,
\bm{r}\cdot\bm{r}\,\FrameBasis
\end{SaveEquation}
where \(\gradS\) acts with respect to the section coordinates \(\bm{r}\). 

\item[Linear strain.]
The first-order truncation of \eqref{eq:exact-gl-strain-Phi} furnishes:
\begin{SaveEquation}{eq:linear-gl-strain-pi}
\hat{\StrainVector}\coloneq
\bm{\gamma}
+\bm{\kappa}\times\bm{r}
+\FrameBasis\otimes\WarpAxialShapes\,\bm{\alpha}'
+\gradS (\WarpAxialShapes\cdot\bm{\alpha})
\end{SaveEquation}
\end{description}

\paragraph{Section kinematic matrix.}
The variation of the strain vector takes the form 
\(\delta\hat{\bm{\varepsilon}} = \TotalStrainMatrix\,\delta\bm{e}\),
where \(\TotalStrainMatrix(\bm{r})\) is the \emph{section kinematic matrix} 
that maps beam strain variations to pointwise strain variations over the section.
From \eqref{eq:wagner-gl-strain-pi}:
\begin{equation}
\TotalStrainMatrix
=\left[
\;\mathbf{1}
~~~-\,\bm{r}^{\times}+\Twist r^2 \FrameBasis\otimes\FrameBasis
~~~~\FrameBasis\otimes\WarpAxialShapes
~~~~\FrameBasis_{\beta}\otimes\nabla_{\beta}\WarpAxialShapes
\right],
\label{eq:def-section-kinematic-matrix-wagner}
\end{equation}
where \(r^2 \coloneq \bm{r}\cdot\bm{r}\).
The linear form \eqref{eq:linear-gl-strain-pi} gives:
\begin{equation}
\TotalStrainMatrix
=\left[
\;\mathbf{1}
~~~-\bm{r}^{\times}
~~~\FrameBasis\otimes\WarpAxialShapes
~~~~\FrameBasis_{\beta}\otimes\nabla_{\beta}\WarpAxialShapes
\right].
\label{eq:def-section-kinematic-matrix-Phi}
\end{equation}

\paragraph{Resultants.}
Recall the Frobenius inner product 
\(
    \bm{A} \cdot \bm{B} \coloneq \operatorname{tr} \bm{A}\bm{B}^{\top}
\)
for two linear operators \(\bm{A}\) and \(\bm{B}\). 
Since \(\delta \hat{\StrainTensor}\) has the representation 
\(\left(\TotalStrainMatrix \delta\FrameStrain\right) \stackrel{s}{\otimes}\FrameBasis\),
\begin{equation*}
    \left<\StressTensor ,\,\delta\StrainTensor\right>_{\Section} = \SectionIntegral{ 
    \operatorname{tr} \left[\StressTensor\, \tfrac{1}{2}\left(\TotalStrainMatrix\delta \FrameStrain \otimes \FrameBasis + \FrameBasis\otimes \TotalStrainMatrix\delta \FrameStrain \right)
    \right]
    }.
\end{equation*}
Exploiting linearity of the trace together with the identity \(\operatorname{tr} (\bm{a}\otimes\bm{b}) = \bm{a}\cdot \bm{b}\) for arbitrary vectors \(\bm{a}\) and \(\bm{b}\) yields
\begin{equation*}
    \left<\StressTensor ,\,\delta\StrainTensor\right>_{\Section} = \delta \FrameStrain\cdot \SectionIntegral{
        \TotalStrainMatrix^{\top}\tfrac{1}{2}\left(\StressTensor \,\FrameBasis + \StressTensor^{\top}\FrameBasis\right)
    },
\end{equation*}
and consequently, requiring \eqref{eq:seSE} section-wise, the conjugate resultants \(\TraceTotalEnergyStress\) are:
\begin{equation}
\TraceTotalEnergyStress = \SectionIntegral{
  \TotalStrainMatrix^{\top} \StressVector
}
\label{eq:stress-resultants}
\end{equation}
where the last equality uses symmetry of \(\StressTensor\) and the definition \eqref{eq:def-stress-vector} of the stress vector.

\iffalse
\subsubsection{Section Constitution}\label{sec:std-section}
The stress condition \eqref{eq:saint-venant-stress} introduces a constraint that determines the plane strains \(\tilde{\StrainTensor}\). This motivates adopting a truncated representation of the strains. Two cases of interest are:

\begin{description}
\item[Moderate (Wagner) strain.]
A truncation of \eqref{eq:exact-gl-strain-Phi} which retains the dominant section-level nonlinearity is:
\begin{SaveEquation}{eq:wagner-gl-strain-pi}
%
\hat{\StrainVector}\coloneq
\bm{\gamma}
+\bm{\kappa}\times\bm{r}
% +\bm{\Phi}\,\bm{\alpha}'
% +\nabla\AxialWarp
+\FrameBasis\otimes\WarpAxialShapes\,\bm{\alpha}'
+\gradS (\WarpAxialShapes\cdot\bm{\alpha})
+\tfrac{1}{2}\Twist^2\,
\bm{r}\cdot\bm{r}\,\FrameBasis
\end{SaveEquation}
where \(\gradS\) acts with respect to the section coordinates \(\bm{r}\). 
From \eqref{eq:wagner-gl-strain-pi}:
\begin{equation}
\TotalStrainMatrix
=\left[
\;\mathbf{1}
~~~-\,\bm{r}^{\times}+\Twist r^2 \FrameBasis\otimes\FrameBasis
~~~~\FrameBasis\otimes\WarpAxialShapes
~~~~\FrameBasis_{\beta}\otimes\nabla_{\beta}\WarpAxialShapes
\right],
\label{eq:def-section-kinematic-matrix-wagner}
\end{equation}
where \(r^2 \coloneq \bm{r}\cdot\bm{r}\).
% Thus, for any virtual increment,
% \begin{equation}
% \delta\hat{\bm{\varepsilon}}
% =\delta\bm{\gamma}
% -\big(\bm{r}
% +\bm{\Phi}\bm{\alpha}\big)\times\delta\bm{\kappa}
% +\bm{\Phi}\,\delta\bm{\alpha}'
% +\bm{\kappa}\times\bm{\Phi}\,\delta\bm{\alpha}
% +\nabla\big(\FrameBasis\cdot\bm{\Phi}\,\delta\bm{\alpha}\big),
% \label{eq:delta-strain-vector-wagner}
% \end{equation}
% and the remainder
% \(
% \delta\tilde{\bm{E}}=\sym\big(\nabla(\mathbf{1}_{\Section}\bm{\Phi}\,\delta\bm{\alpha})\big)
% \)
% is orthogonal to \(\StressTensor\) under the constrained stress condition \eqref{eq:saint-venant-stress} by \cref{prop:sv-orthogonality}.

\item[Linear strain.]
The first-order truncation of \eqref{eq:exact-gl-strain-Phi} furnishes:
\begin{SaveEquation}{eq:linear-gl-strain-pi}
\hat{\StrainVector}\coloneq
\bm{\gamma}
+\bm{\kappa}\times\bm{r}
% +\WarpTotalShapes\,\bm{\alpha}'
+\FrameBasis\otimes\WarpAxialShapes\,\bm{\alpha}'
+\gradS (\WarpAxialShapes\cdot\bm{\alpha})
\end{SaveEquation}

From \eqref{eq:linear-gl-strain-pi}:
\begin{equation}
\TotalStrainMatrix
=\left[
\;\mathbf{1}
~~~-\bm{r}^{\times}
~~~\FrameBasis\otimes\WarpAxialShapes
~~~~\FrameBasis_{\beta}\otimes\nabla_{\beta}\WarpAxialShapes
\right],
\label{eq:def-section-kinematic-matrix-Phi}
\end{equation}
\end{description}

\paragraph{Resultants.}
Recall the Frobenius inner product 
\(
    \bm{A} \cdot \bm{B} \coloneq \operatorname{tr} \bm{A}\bm{B}^{\top}
\)
for two linear operators \(\bm{A}\) and \(\bm{B}\). 
Since \(\delta \hat{\StrainTensor}\) has the representation \(\left(\TotalStrainMatrix \delta\FrameStrain\right) \stackrel{s}{\otimes}\FrameBasis\) (modulo the in-plane term which is orthogonal to \(\StressTensor\)), 
\begin{equation*}
    \left<\StressTensor ,\,\delta\StrainTensor\right>_{\Section} = \SectionIntegral{ 
    \operatorname{tr} \left[\StressTensor\, \tfrac{1}{2}\left(\TotalStrainMatrix\delta \FrameStrain \otimes \FrameBasis + \FrameBasis\otimes \TotalStrainMatrix\delta \FrameStrain \right)
    \right]
    }.
\end{equation*}
Exploiting linearity of the trace together with the identity \(\operatorname{tr} (\bm{a}\otimes\bm{b}) = \bm{a}\cdot \bm{b}\) for arbitrary vectors \(\bm{a}\) and \(\bm{b}\) yields
\begin{equation*}
    \left<\StressTensor ,\,\delta\StrainTensor\right>_{\Section} = \delta \FrameStrain\cdot \SectionIntegral{
        \TotalStrainMatrix^{\top}\tfrac{1}{2}\left(\StressTensor \,\FrameBasis + \StressTensor^{\top}\FrameBasis\right)
    },
\end{equation*}
and consequently, requiring \eqref{eq:seSE} section-wise, the conjugate resultants \(\TraceTotalEnergyStress\) are:
\begin{equation}
\TraceTotalEnergyStress = \SectionIntegral{
  \TotalStrainMatrix^{\top} \StressVector
}
\label{eq:stress-resultants}
\end{equation}
where the last equality uses symmetry of \(\StressTensor\) and the definition \eqref{eq:def-stress-vector} of the stress vector.
\fi

\subsubsection{Summary}

\begin{ambox}[\gls{frame:exact} \gls{frame:shear} Frame (SVQ)]{box:exact-shear-svq}

\textbf{Equilibrium} \eqref{eq:finite-equilibrium}

\begin{equation*} 
\left\{\begin{array}{r}
(\FrameRotation \bm{n})^{\prime}+{\bm{f}}_n=\mathbf{0} \\
(\FrameRotation \bm{m})^{\prime}+\bm{x}^{\prime} \times (\FrameRotation \bm{n}) +{\bm{f}}_m=\mathbf{0} \\
\bm{w}^{\prime}-\bm{v} + \bm{f}_{\rm w} = \mathbf{0}
\end{array}\right. 
\end{equation*}

\textbf{Kinematics} \eqref{eq:finite-alignment-field}
\begin{equation*} 
  \begin{aligned}
    \ShearStrain &= \FrameRotation^{\top}\bm{x}^\prime - \FrameBasis 
  \end{aligned}
  \qquad\text{ and }\qquad
  \begin{aligned}
    \CurveStrain &= 
    \left(\FrameRotation^{\top}\FrameRotation^\prime\right)^{\vee}
  \end{aligned}
\end{equation*}

\end{ambox}


\begin{ambox}[Basic Shear Frame]{box:basic-shear}

\textbf{Equilibrium}

\begin{equation}
-\EquilibriumDerivative\bm{A}_* \bm{s}_{\mathrm{p}} = \left\{\begin{aligned}
& \ShearResult^{\prime} \\
& \CurveResult^{\prime} + \FrameBasis \times\ShearResult \\
& \bm{w}' - \bm{v}
\end{aligned}
\right.
\end{equation}

\textbf{Kinematics.} \eqref{eq:gam-kap-linear}
$$
\begin{aligned}
\ShearStrain &= \bm{u}^{\prime}+\FrameBasis  \times \bm{\theta} \\
\CurveStrain &= \bm{\theta}^{\prime}
\end{aligned}
$$

\textbf{Operators}
\[
  \mathbb{B}
  =
  -\begin{bmatrix}
    \mathbf{1}\,\partial_x & 0 & 0 & 0 \\
    \FrameBasis^{\times} & \mathbf{1}\,\partial_x & 0 & 0 \\
    0 & 0 & \mathbf{1}\,\partial_x & -\mathbf{1}
  \end{bmatrix},
  \qquad
  \mathbb{B}^*
  =
  \left[\begin{array}{cc|c}
    \mathbf{1}\,\partial_x & \FrameBasis^{\times} & 0 \\
    0 & \mathbf{1}\,\partial_x & 0 \\[0.3ex]
    \hline
    0 & 0 & \mathbf{1}\,\partial_x \\
    0 & 0 & \mathbf{1}
  \end{array}\right].
\]
\end{ambox}
